\section{问题背景与重述}

低空经济中的基础设施巡检具有点位分散、任务重复、时间限制严格等特点。多架无人机从统一基地同时起飞，既要完成不同等级巡检点规定的有效到达次数，又要在最大工作时间内返回基地。若进一步考虑按时间启用的圆形禁飞区，路线可行性不再只由空间距离决定，而由飞行位置和到达时刻共同决定。

题目给出四个 Case 的巡检点坐标、巡检等级以及动态禁飞区参数。所有无人机同型号，平均巡航速度为 $v=55\ \mathrm{km/h}$，单次巡检耗时 $s=5/60=1/12\ \mathrm{h}$，附件坐标单位对应 $0.1\ \mathrm{km}$，基地坐标为 $(0,0)$，每天 8:00 同时出发。问题一要求在单架无人机工作时间不超过 $H=9\ \mathrm{h}$ 的前提下，先确定最低可行机队，再优化最大完成时间；问题二要求在问题一机队基础上改善各机工作时间均衡；问题三加入带启停时间的圆形禁飞区，重新规划动态可行路线。

本文的统一主线是：先将原始点位和巡检等级转化为有效到达任务，再以问题一得到的机队和路线为基础逐步加入负载均衡与时空禁飞约束。图~\ref{fig:roadmap} 展示了三个问题的继承关系。

\clearpage
\begin{landscape}
\begin{figure}[H]
\centering
\includegraphics[width=0.94\linewidth]{fig_roadmap.pdf}
\caption{三个子问题的统一技术路线}
\label{fig:roadmap}
\end{figure}
\end{landscape}
\clearpage

\subsection{问题一：最低可行机队与最大完成时间}

对每个物理点，根据巡检等级展开为若干带唯一 \texttt{task\_id} 的任务。任务可由不同无人机分担，但同一条路线不允许相邻任务属于相同物理点。模型先计算理论机队下界，从下界开始检查固定机队是否能完成全部任务、闭环返回基地且满足 $9\ \mathrm{h}$ 限制；首次通过独立验证的数量记为最低已验证可行数量。若该数量下仍有优化空间，再以 $C=\max_k T_k$ 为目标压缩最后完成时间。

\subsection{问题二：固定机队的工作均衡}

问题二不重新定义任务，而是固定问题一得到的 $N_{\min}$ 和全部硬约束。设 $L=\min_kT_k$、$\delta=C-L$，采用字典序目标先保持 $C$ 不劣于问题一基准值，再最小化 $\delta$，最后以总距离 $D$ 区分同等级解。这样可以避免为两个没有题面权重的目标人为指定加权系数。

\subsection{问题三：动态禁飞区下的时空重规划}

问题三将静态边距离替换为依赖出发时刻的最早安全到达函数。对一条边先做线段与圆的几何判断，再判断实际进入和离开禁飞圆的时段是否与 $[t_s,t_e)$ 重叠。发生冲突时比较安全等待、切线—圆弧绕行和调整任务顺序；每次等待或绕行都必须重新传播后续时刻，并复核巡检服务区间、基地停留和 $9\ \mathrm{h}$ 上限。若固定问题一机队无法恢复可行性，只逐架增加无人机。
